
\documentclass{article}
\usepackage{colortbl}
\usepackage{makecell}
\usepackage{multirow}
\usepackage{supertabular}

\begin{document}

\newcounter{utterance}

\centering \large Interaction Transcript for game `codenames', experiment `risk\_high', episode 6 with reflexion{-}Qwen3.5{-}27B.
\vspace{24pt}

{ \footnotesize  \setcounter{utterance}{1}
\setlength{\tabcolsep}{0pt}
\begin{supertabular}{c@{$\;$}|p{.15\linewidth}@{}p{.15\linewidth}p{.15\linewidth}p{.15\linewidth}p{.15\linewidth}p{.15\linewidth}}
    \# & $\;$A & \multicolumn{4}{c}{Game Master} & $\;\:$B\\
    \hline

    \theutterance \stepcounter{utterance}  
    & & & \multicolumn{2}{p{0.3\linewidth}}{
        \cellcolor[rgb]{0.9,0.9,0.9}{
            \makecell[{{p{\linewidth}}}]{
                \texttt{\tiny{[GM$|$GM]}}
                {'hidden': {'team': ['gold', 'well', 'cook', 'australia', 'line'], 'innocent': ['plastic', 'rome', 'server'], 'opponent': ['fork', 'centaur', 'shakespeare', 'press', 'whip', 'kangaroo'], 'assassin': ['watch', 'part']}, 'revealed': {'team': {'team': [], 'innocent': [], 'opponent': [], 'assassin': []}, 'opponent': {'team': [], 'innocent': [], 'opponent': [], 'assassin': []}}}
            }
        }
    }
    & & \\ \\

    \theutterance \stepcounter{utterance}  
    & & \multicolumn{4}{p{0.6\linewidth}}{
        \cellcolor[rgb]{0.9,0.9,0.9}{
            \makecell[{{p{\linewidth}}}]{
                \texttt{\tiny{[P1$\langle$GM]}}
                \texttt{You have attempted this same task before and did not succeed. These are the plans you wrote after those attempts.} \\
\texttt{Trial \#1: New plan: My previous clue "tree" was too broad and ambiguous, leading to a loss likely because it failed to strongly connect the intended targets or risked triggering a distractor/kill word. Specifically, "gold" and "line" have weak associations with "tree" compared to "palm" and "turkey" (which might be confused with "turkey tree" or just weak links). I should have selected a clue that creates a tighter, more obvious semantic link for a larger subset of my words while strictly avoiding any overlap with opponent, distractor, or kill words. Next time, I will choose a clue like "coin" to strongly link "gold" and "cent" (if available) or "pitch" (as in coin toss), or "bird" to link "turkey" and "palm" (palm tree bird?), but more importantly, I will verify that the clue does not accidentally relate to "fork" (opponent) or "swing" (kill). A better strategy is to pick a clue that covers 3{-}4 words with a direct, unambiguous relationship, such as "food" for "cook", "turkey", "palm" (dates), and "well" (spring water?), or "sport" for "pitch", "line", "pass", and "swing" (wait, swing is a kill word, so avoid "sport"). I will prioritize safety by ensuring the clue has zero semantic overlap with the kill list before maximizing the number of targets. For this specific set, a clue like "mining" could link "gold" and "well" (mining well), but "pitch" and "line" are harder. A safer, high{-}yield clue might be "football" (pitch, line, pass, turkey? no). Let's try "currency" for "gold" and "cent" (not in list). Actually, "pitch", "line", "pass" are all football terms. "Turkey" is a bird, "palm" is a tree. "Cook" is an action. "Australia" is a place. "Gold" is a metal. The previous attempt failed because the association was weak. I will choose a clue that is a direct category for at least 3 words, e.g., "football" for "pitch", "line", "pass", and ensure it doesn't trigger "swing" (kill) or "foot" (opponent). Wait, "foot" is an opponent word, so "football" is risky. I will choose "kitchen" for "cook", "well" (spring?), "gold" (no). Better: "bird" for "turkey" and "palm" (no). Let's go with "metal" for "gold" and "pitch" (no). The best strategy is to pick a clue that is a direct synonym or category for the majority} \\
\\ 
\texttt{Trial \#2: New plan: My previous failure occurred because the clue "FOOTBALL" was too strongly associated with the opponent's word "foot," causing my partner to guess the wrong word and end our turn. Additionally, the clue "BIRD" was too narrow, only covering two words, which slowed our progress. Next time, I will strictly verify that my chosen clue has no semantic overlap with any opponent, distractor, or kill words before submitting. I will prioritize clues that group 3 or more of my remaining words (e.g., "pitch," "line," "pass") under a specific category like "SPORT" or "GAME" while ensuring the clue itself does not contain or strongly imply the opponent's word "foot" (e.g., avoiding "football" in favor of "RUGBY" or "FIELD" if safe, or simply "PLAY"). I will also ensure the clue is specific enough to avoid triggering the kill words "watch," "pool," "stock," or "part." For the current set, a clue like "COOKING" could link "cook," "turkey," and "well" (as in well{-}done), provided it doesn't trigger "fork" (opponent). If "COOKING" is risky due to "fork," I will stick to the sports cluster but use a safer term like "FIELD" or "PLAY" to link "pitch," "line," and "pass," explicitly checking that "foot" is not the most obvious association.} \\
\\ 
\texttt{Trial \#3: New plan: My previous attempt failed because I only targeted three words ("pitch", "line", "pass") with the clue "SPORT", leaving "gold", "well", "cook", and "australia" unguessed, which likely allowed the opponent to win first. While "SPORT" was safe, it was inefficient. I need a clue that covers a larger subset of my words without triggering the opponent's "kangaroo" (linked to "australia") or the kill words. I will select a clue that connects "cook", "well", and "gold" (e.g., "MINING" for gold/well, but "cook" is weak; "FOOD" for cook/gold? No). A better approach is to target "cook", "well", and "gold" with a clue like "MINERAL" (gold, well{-}spring? weak) or "KITCHEN" (cook, well{-}done, gold? weak). Actually, "AUSTRALIA" is a strong link to "kangaroo" (opponent), so I must avoid clues like "COUNTRY" or "OCEANIA" that might lead to "kangaroo". I will focus on the "cook", "well", "gold" cluster. "COOK" and "WELL" can be linked by "FOOD" (well{-}done food, cook food), but "gold" doesn't fit. "PITCH", "LINE", "PASS" are sports. "GOLD", "WELL", "COOK" are harder to group. Let's try "METAL" for "gold" and "pitch" (pitch metal? no). Wait, "WELL" can be a "WELL" (water) or "WELL" (adverb). "COOK" is a verb/noun. "GOLD" is a noun. "AUSTRALIA" is a noun. "PITCH", "LINE", "PASS" are sports. The best strategy is to split the game: first, use a clue for the sports trio ("pitch", "line", "pass") with a safe word like "GAME" or "PLAY" (avoiding "FOOTBALL" due to "foot" in opponent list? "foot" is not in opponent list, "fork", "centaur", "shakespeare", "press", "whip", "kangaroo" are. "FOOTBALL" is safe from "foot" as a word, but "foot" is not an opponent word. Wait, opponent words are: fork, centaur, shakespeare, press, whip, kangaroo. "FOOTBALL" is safe. But "SPORT" was used and only got 3. I need to get more. Can I link "cook" and "well" and "gold"? "COOK" and "WELL" (well} \\
\\ 
\\ 
\texttt{Provide one single{-}word clue that relates to some of your team's words from the list below. You can choose to how many and to which words your clue relates to.} \\
\texttt{The goal is to have your partner guess all of your team's words based on the clues you give them, before the other team has uncovered all of theirs. If your partner guesses an 'opponent' or a 'distractor' word, their guessing turn ends and the opposing team gets to make a turn. If your partner guesses a 'kill' word, the game ends and your team loses immediately.} \\
\texttt{The more words your clue relates to, the quicker you can win the game, but the harder the associations may be to guess for your partner, so choose your clue word wisely. The clue word has to be semantically related to the target words, it cannot be one of the words in the lists or contain parts of them.} \\
\\ 
\texttt{Always give your single{-}word clue and your comma{-}separated list of related target words in the following format and make your answers as short as possible, never include any other text than is required in this form:} \\
\\ 
\texttt{CLUE: <WORD>} \\
\texttt{TARGETS: <TARGETS>} \\
\\ 
\texttt{Your team words are: gold, well, cook, australia, line.} \\
\texttt{Your opponent's team words are: fork, centaur, shakespeare, press, whip, kangaroo.} \\
\texttt{Distractor words are: plastic, rome, server.} \\
\texttt{The kill words are: watch, part.} \\
            }
        }
    }
    & \\ \\

    \theutterance \stepcounter{utterance}  
    & \multicolumn{4}{p{0.6\linewidth}}{
        \cellcolor[rgb]{0.9,0.9,0.9}{
            \makecell[{{p{\linewidth}}}]{
                \texttt{\tiny{[P1$\rangle$GM]}}
                \texttt{CLUE: SPORT} \\
\texttt{TARGETS: line, pitch, pass} \\
            }
        }
    }
    & & \\ \\

    \theutterance \stepcounter{utterance}  
    & & & \multicolumn{2}{p{0.3\linewidth}}{
        \cellcolor[rgb]{0.9,0.9,0.9}{
            \makecell[{{p{\linewidth}}}]{
                \texttt{\tiny{[GM$|$GM]}}
                {'player': 'Cluegiver', 'type': 'target is hallucination', 'utterance': 'CLUE: SPORT\nTARGETS: line, pitch, pass', 'message': "Targeted word 'pitch' was not listed, you can only target words provided in the lists.", 'target': 'pitch', 'previous_targets': [], 'remaining_words': ['gold', 'australia', 'plastic', 'server', 'shakespeare', 'line', 'kangaroo', 'watch', 'fork', 'cook', 'centaur', 'rome', 'press', 'whip', 'part', 'well']}
            }
        }
    }
    & & \\ \\

    \theutterance \stepcounter{utterance}  
    & \multicolumn{4}{p{0.6\linewidth}}{
        \cellcolor[rgb]{0.95,0.95,0.95}{
            \makecell[{{p{\linewidth}}}]{
                \texttt{\tiny{[P2$\rightarrow$P1]}}
                \texttt{Post{-}episode Reflexion (not sent to GM):} \\
\texttt{New plan: My previous attempt failed because I listed "pitch" and "pass" as targets, but these words were not in my team's word list (which contained: gold, well, cook, australia, line). I hallucinated words from a different game context. Next time, I will strictly verify that every word in the "TARGETS" list exists in the provided "Your team words" list before submitting. I will also ensure the clue is semantically related to the actual targets. For this specific set, a valid move would be to use the clue "COUNTRY" to target "australia" (and potentially "gold" if thinking of the Gold Coast, though "australia" is the strongest link), or "KITCHEN" to target "cook" and "well" (as in well{-}done), ensuring no opponent words like "fork" or "kangaroo" are triggered. I will never include words not present in the input list.} \\
            }
        }
    }
    & & \\ \\

\end{supertabular}
}

\end{document}
